% !TeX root = exam_template.tex
% ==========================================
% 湖北大学课程考试试题纸 - 试卷题目内容
% ==========================================

\examsection{一、选择题（每题 2 分，共 20 分）}
\begin{questions}
  \item 设正实数 $x, y, z$ 满足关系式 $x+y+z=1$，则函数 $f(x,y,z) = \frac{1}{x} + \frac{4}{y} + \frac{9}{z}$ 的最小值为\answerline[15mm]。
    \choice{14}{25}{36}{49}
  \item 一个四位数的质数 $p$。当 $p$ 被 30 除时，可能的余数个数为\answerline[15mm]。
    \choice{8}{10}{15}{30}
  \item 7 位成员在进行舞台排练时，要求 3 位主唱两两均不相邻。则这 7 位成员排成一排的队形种数为\answerline[15mm]。
    \choice{720}{1440}{2880}{5040}
  \item 蒙德、璃月、稻妻三个国家之间的距离组成三角形的三边 $a, b, c$，且满足关系式 $a^3 + b^3 = c^3$。则关于该三角形的内角 $C$，下列说法正确的是\answerline[15mm]。
    \choice{$C > 90^\circ$}{$C < 90^\circ$}{$C = 90^\circ$}{无法确定}
  \item 函数 $f: \mathbb{R} \to \mathbb{R}$ 满足方程 $f(x+y) = f(x)f(y) - f(xy) + 1$。若已知 $f(1) = 2$，则 $f(2026)$ 的值为\answerline[15mm]。
    \choice{2025}{2026}{2027}{1}
  \item 四次多项式 $P(x)$。若 $P(1) = 1, P(2) = 2, P(3) = 3, P(4) = 4$，则 $P(5)$ 为\answerline[15mm]。
    \choice{5}{24}{29}{125}
  \item 序列 $\{a_n\}$ 满足首项 $a_1 = 1$，且递推式为 $a_{n+1} = a_n + \frac{1}{a_n}$ ($n \ge 1$)。则关于项 $a_{100}$ 的估算，正确的是\answerline[15mm]。
    \choice{$10 < a_{100} < 15$}{$15 < a_{100} < 20$}{$20 < a_{100} < 25$}{$a_{100} > 50$}
  \item 一个平面结界战场中，分布着 10 个战斗节点（任意三点不共线）。由这些节点两两连成的所有线段中，在凸包内部的交点个数最多为\answerline[15mm]。
    \choice{45}{120}{210}{252}
  \item  $N = 2^{10} \times 3^5$。则 $N$ 的正约数中，是 6 的倍数的约数共有\answerline[15mm] 个。
    \choice{11}{50}{66}{72}
  \item 一个魔法阵的 7 个顶点在复平面上代表方程 $z^7 - 1 = 0$ 的非 1 复数根 $\omega, \omega^2, \dots, \omega^6$。则代数式 $(1-\omega)(1-\omega^2)\cdots(1-\omega^6)$ 的值为\answerline[15mm]。
    \choice{0}{1}{7}{49}
\end{questions}

\examsection{二、填空题（每空 1 分，共 10 分）}
\begin{questions}
  \item 对于正实数 $a, b$，若满足 $a+b=4$，则当 $a = $ \fillblank 且 $b = $ \fillblank 时，代数式 $a^2b$ 取得最大值。
  \item 正整数密钥 $N$ 满足 $N \equiv 3 \pmod 5$ 且 $N \equiv 4 \pmod 7$，则 $N$ 的最小正整数解为 \fillblank；若在此基础上还要求 $N \equiv 2 \pmod{11}$，则其最小正整数解为 \fillblank。
  \item 露天舞台的四个机位 $A, B, C, D$ 在平面上依次共圆。已知测得距离分别为 $AB = 3, BC = 4, CD = 5, DA = 6$。则四边形对角线乘积 $AC \times BD =$ \fillblank，该四边形的面积为 \fillblank。
  \item 动漫社将 10 个相同的周边分配给 3 位核心社员。若要求每位社员至少分配到 1 个，则分配方案数有 \fillblank 种；若要求社员甲分配到的周边个数必须为偶数，则分配方案数有 \fillblank 种。
  \item 三位乐评人 A, B, C 投票给一首歌曲投票的概率分别为 $\frac{1}{2}, \frac{2}{3}, \frac{3}{4}$，且三人投票相互独立。则该歌曲恰好获得 2 票的概率为 \fillblank，至少获得 1 票的概率为 \fillblank。
\end{questions}

\examsection{三、简答题（共 10 分）}
\begin{questions}
  \item （6 分）在某次有 6 位明星参加的晚宴中，任意两位明星之间要么是互相关注的粉丝关系，要么是互不认识的关系。请利用图论或组合数学方法证明：在他们之中，必然存在 3 位明星，这 3 人两两之间要么都是互相关注的关系，要么都是互不认识的关系。
    \answerspace{45mm}
  \item （4 分）设 $x_1, x_2, x_3$ 是关于 $x$ 的三次方程 $x^3 - 3x^2 + x + 1 = 0$ 的三个根。求代数式 $x_1^2 + x_2^2 + x_3^2$ 的值，并写出求解步骤。
    \answerspace{35mm}
\end{questions}

\examsection{四、计算题（共 60 分）}
\begin{questions}
  \item （10 分）炼成能量转换中，魔力输出受到三个正实数因子 $x, y, z$ 的制约，且满足约束条件：$x^2 + 2y^2 + 3z^2 = 6$。
  请利用柯西-施瓦茨不等式（Cauchy-Schwarz Inequality），求魔力合成函数 $f(x,y,z) = x + 2y + 3z$ 的最大值，并指出取得最大值时 $x, y, z$ 的值。
    \answerspace{48mm}
  \item （10 分）为了破译世界线变动率的安全密钥，需要求出以下巨大整数的模：
  $$N = 3^{2026} + 4^{2026} \pmod{13}$$
  请利用费马小定理（Fermat's Little Theorem）计算并求出 $N$ 在模 13 下的最小非负剩余。
    \answerspace{48mm}
  \item （15 分）一级魔法使考试的结界呈三角形 $ABC$。已知点 $D$ 在边 $BC$ 上，且满足 $BD : DC = 2 : 1$。点 $E$ 在边 $AC$ 上，满足 $AE : EC = 3 : 4$。线段 $AD$ 与 $BE$ 相交于点 $P$，延长线段 $CP$ 交 $AB$ 于点 $F$。\\
  (1) 利用塞瓦定理（Ceva's Theorem）求出 $AF : FB$ 的比值；（5 分）\\
    \answerspace{50mm}
  (2) 利用梅涅劳斯定理（Menelaus's Theorem）求出 $AP : PD$ 的比值；（5 分）\\
    \answerspace{80mm}
  (3) 求三角形 $ABP$ 的面积与三角形 $ABC$ 总面积的比值 $\frac{S_{\triangle ABP}}{S_{\triangle ABC}}$。（5 分）
    \answerspace{70mm}
  \item （13 分）颁奖典礼的后台，有 $n$ 名歌手。他们每人带了一份准备送给其他人的礼物。现在进行随机交换，要求每人正好得到一份别人的礼物，且没有任何人得到自己带来的礼物（全错位排列）。记 $D_n$ 为 $n$ 个元素的全错位排列数。\\
  (1) 证明递推公式：$D_n = (n-1)(D_{n-1} + D_{n-2})$ 对 $n \ge 3$ 成立，其中 $D_1 = 0, D_2 = 1$；\\（6 分）
    \answerspace{70mm}
  (2) 设 $n = 5$（即有 5 名歌手进行礼物交换），求所有可能的礼物交换方案数 $D_5$；\\（3 分）
    \answerspace{70mm}
  (3) 若这 5 名歌手随机抽取礼物，求恰好有 2 个人抽到自己礼物的概率。（4 分）
    \answerspace{70mm}
  \item （12 分）设 $\omega = \cos\frac{2\pi}{11} + i \sin\frac{2\pi}{11}$ 为 11 次单位根。\\
  (1) 证明：$\sum_{k=1}^{10} \omega^k = -1$；（4 分）
  (2) 计算下列三角函数和的值：
      $$S = \cos\frac{2\pi}{11} + \cos\frac{4\pi}{11} + \cos\frac{6\pi}{11} + \cos\frac{8\pi}{11} + \cos\frac{10\pi}{11}$$
      （写出详细的计算步骤）。（8 分）
    \answerspace{58mm}
\end{questions}
